\section{Tamamlayıcı vaka: bloklu bir tarla deneyi}
\label{sec:v1-colt-park}

\textbf{Soru ve kaynak.} Colt Park deneyinde çiftlik gübresi (FYM) ile
gübresiz kontrolün kuru ot verimleri nasıl karşılaştırılır?
\citet{villagalaviz2023} çalışmasının S1 Data ve S1 File ekleri kullanılır.
Bu, makalenin çok göstergeli analizinin tekrarı değil, tek yanıtlı bir
öğretim uyarlamasıdır. Karşılaştırma ve yanıt bu vaka için belirlenmiştir;
yayımlanmış sonuçlardan habersiz yapılmış bir önkayıt değildir.

\textbf{Atama ve kayıt kapsamı.} Yöntem ekinin 2. sayfasına göre 72
parselin tohum eklenmeyen 24'ü incelenir. Her biri 15 m$^2$ olan
parseller üç blokta yer alır. Her blokta FYM, NPK, NPK+FYM ve gübresiz
kontrol için ikişer parsel bulunur; gübreleme rejimleri blok içinde
rastgele atanmıştır. Veri tablosunun blok--rejim sayımları bu açıklamayla
eşleşir. Burada yalnız FYM ve kontrol seçilerek her bloktan dört,
toplam 12 parsel kullanılır. Diğer rejimler ham dosyadan silinmez.

\begin{researchbox}
\textbf{Ölçüm birimi ve izlenebilirlik.} Yanıt, 2011--2014 ortalama kuru
madde verimidir; yıllar ayrı bağımsız parseller değildir. Excel sözlüğü
alanı m$^2$, yöntem ekinin hesap açıklaması ağırlığı gram olarak belirtir;
burada birim g/m$^2$ olarak korunur. Bu değer 15 m$^2$ parselin toplam
ağırlığı veya 2017'deki besin içeriği ölçümü değildir.

Ek dosyada fiziksel parsel numarası ve yıllık ham ölçümler yoktur.
Üretilen \texttt{source\_row}, Excel'in DATA sayfasındaki satırdır;
saha kimliği değildir. Analize giren 12 kayıtta eksik değer yoktur;
bu kontrol, geçmişte hiç ölçüm kaybı olmadığını kanıtlamaz. Eski Spector
örneğinden farklı olarak blok içi atama burada yöntem ekiyle desteklenir;
ancak özgün saha atama tutanağı yeniden denetlenmiş değildir.
\end{researchbox}

\textbf{Analiz tanımı.} Farkın yönü FYM eksi kontroldür. Her blokta iki
grup ortalaması çıkarılır; üç blok farkına eşit ağırlık verilir:
\[
\widehat{\tau}=\frac{1}{3}\sum_{h=1}^{3}
\bigl(\overline{Y}_{h,\mathrm{FYM}}-\overline{Y}_{h,\mathrm{kontrol}}\bigr).
\]
Her blokta aynı sayıda parsel bulunduğu için bu tahmin, iki grubun genel
ortalamalarının farkına da eşittir. Eşleştirilmiş öğrenci çiftleri gibi
parselleri satır sırasıyla eşleştirmek veya blokları yok saymak gerekmez.

\begin{table}[H]
\centering
\caption{Colt Park: blok ortalamaları ve FYM--kontrol farkı (g/m$^2$)}
\label{tab:v1-colt-park}
\begin{tabular}{@{}lrrr@{}}
\toprule
Blok & Kontrol & FYM & Fark \\
\midrule
ONE & 467,563 & 619,629 & 152,066 \\
TWO & 372,961 & 548,338 & 175,377 \\
THREE & 376,764 & 563,784 & 187,020 \\
\midrule
Eşit ağırlıklı ortalama & 405,763 & 577,250 & 171,488 \\
\bottomrule
\end{tabular}
\end{table}

\textbf{Tasarıma dayalı test.} Keskin sıfır hipotezi, seçilen iki rejimin
ilgili parsellerin hiçbirinin yanıtını değiştirmemesidir. NPK ve NPK+FYM
parsellerinin yerleri sabit tutulduğunda, kalan dört parselin ikisine FYM
etiketi vermenin her blokta $\binom{4}{2}=6$ yolu vardır. Bloklar birlikte
$6^3=216$ eş olasılıklı koşullu atama verir. Bunların tamamında
$\widehat{\tau}$ yeniden hesaplanır. İki yönlü uçluk
$|\widehat{\tau}_{\rm atama}|\geq|\widehat{\tau}_{\rm gözlenen}|$
ile tanımlanır. İki atama bu koşulu sağlar:
\[
p_{\rm tam}=\frac{2}{216}=0{,}009259\ldots
\]
Bu, Monte Carlo benzetimi değildir; tohum gerekmez ve pay/paydaya bir
eklenmez. Test, belgelenen atama düzeninin eş olasılıklı koşullu
dağılımına dayanır; farklı bir atama kuralı için aynı dağılım kullanılamaz.

\textbf{Belirsizliği ayrıca belirtme.} Tam test bir güven aralığı değildir.
Tamamlayıcı olarak üç sabit blok etkisi ve ortak FYM etkisi içeren
$Y_{hi}=\alpha_h+\tau Z_{hi}+\varepsilon_{hi}$ modeli kurulmuştur;
$Z=1$ FYM, $Z=0$ kontroldür. Bağımsız, normal, ortak varyanslı hatalar ve
bloklar arasında ortak eklemeli etki varsayımı altında klasik yüzde 95
$t$ aralığı $[106{,}546;\ 236{,}429]$ g/m$^2$'dir. Standart hata 28,162;
artık serbestlik derecesi $12-4=8$'dir. Bu model aralığı tam
rastgeleleştirme testinin tersinden elde edilmemiştir; varsayımları yalnız
küçük $p$ değeriyle doğrulanmaz. Üç blok ve 12 parselde varsayım
incelemesinin gücü sınırlıdır.

\textbf{Raporlama örneği.} ``FYM ve kontrol için altışar parselin
2011--2014 ortalama kuru madde verimleri karşılaştırılmıştır.
Bloklara eşit ağırlıklı fark 171,488 g/m$^2$; blok içi tam
rastgeleleştirme testinde iki yönlü $p=0{,}009259$ bulunmuştur.
Eklemeli blok modelinin klasik yüzde 95 güven aralığı
$[106{,}546;\ 236{,}429]$ g/m$^2$'dir; bu aralık ek model
varsayımlarına dayanır. Rastgele atama, uygun uygulama ve parseller arası
etkileşimsizlik koşulları altında bu düzen içindeki nedensel karşılaştırmayı
destekler. Sonuç tüm çayırlar için bir gübre reçetesi, biyolojik çeşitlilik
sonucu veya kısa süreli tek gübre uygulamasının etkisi değildir.''

\textbf{Çalıştırma ve doğrulama.} Ham ekler, veri sözlüğü, dönüşüm,
216 atama ve sayısal çıktılar
\nolinkurl{companion/vakalar/v1-colt-park/} klasöründedir. Paket kökünde:
\begin{lstlisting}[style=prismPythonApp,language={}]
python -B companion/vakalar/v1-colt-park/analyze.py
\end{lstlisting}
Yüklenen iki Excel aynı dosyanın kopyasıdır; iki bağımsız veri kaynağı
sayılmaz. Yazarların R kodu yüklenmediği için özgün R analizi çalıştırılmadı;
buradaki Python sonuçları ham ekteki tek yanıt üzerinden hesaplandı.

\textbf{Görevler ve gerekçeli yanıtlar.}
\begin{enumerate}[label=\arabic*.]
\item \emph{Neden 48 gözlem değil 12 parsel?} Dört yılın ortalaması
parsel başına tek yanıttır; yıllık ham kayıtlar dosyada yoktur.
\item \emph{Neden $\binom{12}{6}$ değil $6^3$ atama?} Blok başına iki
FYM ve iki kontrol korunur; bloklar arasında serbest etiket değişimi
başka bir atama düzenini temsil eder.
\item \emph{Satır sırasına göre altı çift kurabilir miyiz?} Hayır;
belgelenmiş yapı üç bloktur, Excel sırası deneysel eşleştirme değildir.
\item \emph{Tam $p$ değeri ile model aralığı aynı dayanağa mı sahip?}
Hayır; aralık ortak eklemeli etki ve hata dağılımı varsayımlarını ekler.
\item \emph{Spector'dan fark ve genelleme sınırı nedir?} Burada atama
belgesi vardır; buna rağmen rastgele atama, sahanın bütün çayırlar
arasından rastgele seçildiğini göstermez.
\end{enumerate}
\textbf{Değerlendirme (10 puan).} Her görevde doğru ayrım 1 puan,
bu vakaya özgü gerekçe 1 puandır. Yalnız doğru sayıyı yazmak gerekçe
puanını getirmez. Nedensellik ile genellenebilirliği eş tutan son yanıtın
gerekçe puanı verilmez. Bu ölçütler öğrenciyle yapılmış bir etkililik
araştırmasının sonucu değildir.